\begin{resumo}[Abstract]
\begin{otherlanguage*}{english}
This document formalizes the Software Architecture Specification for SIGEA-GTT (Intelligent Management System for Adverse Events -- Global Trigger Tool), a hospital computational platform dedicated to retrospective record review, clinical training, and healthcare epidemiological surveillance, developed through an interdisciplinary partnership between the Department of Computing (DCOMP) and Department of Nursing at the Federal University of Sergipe (UFS). Structured according to Philippe Kruchten's renowned 4+1 Architectural View Model, this work characterizes the system through five complementary perspectives: (i) the \textbf{Scenario View (+1)}, anchoring architectural decisions to critical use cases including restricted institutional authentication, retrospective auditing under a strict 20-minute countdown timer, seamless integration of six Quality Management Tools, and automated calculation of IHI hospital epidemiological indicators ($TI_{1000}$, $TI_{100}$, $P_{EA}$); (ii) the \textbf{Logical View}, detailing the layered architecture grounded in Clean Architecture and Domain-Driven Design (DDD) within the Spring Boot 3.3 (Java 21) backend, alongside the reactive standalone component and Signals paradigm in Angular 18, accompanied by domain class and entity-relationship diagrams (ERD) in TikZ; (iii) the \textbf{Process View}, addressing concurrency, HikariCP connection pooling, ACID transaction management, reactive client-side timer life cycles, and UML sequence diagrams in TikZ for core operational workflows; (iv) the \textbf{Development View}, outlining monorepo organization, Maven and Angular package structures, a software component diagram, and the strict Quality Gate mandating 100\% automated test coverage and 100\% codebase documentation; and (v) the \textbf{Physical / Deployment View}, defining node topology, container orchestration, and HTTPS/TLS communication channels, while enforcing absolute data isolation from University Hospital (HU-UFS/EBSERH) production environments in strict compliance with the Brazilian General Data Protection Law (LGPD). Furthermore, the document provides an architectural traceability matrix aligning Kruchten's views with the 10 Non-Functional Requirements of the ISO/IEC 25010 standard.

\vspace{\onelineskip}
\noindent
\textbf{Keywords}: Software Architecture. Kruchten 4+1 View Model. Global Trigger Tool. Clean Architecture. Patient Safety. Federal University of Sergipe.
\end{otherlanguage*}
\end{resumo}
